\section{Lean formalization architecture}
\label{sec:lean}

\subsection{Existing interfaces}

The implementation should extend, not duplicate, the following inspected
modules.

\begin{longtable}{@{}p{0.37\textwidth}p{0.55\textwidth}@{}}
\toprule
Module/declaration & Role \\
\midrule
\endhead
\code{ConciseIPRA/PR.lean}: \code{PR}, \code{PTm} & Fixed primitive-recursive
function symbols and arithmetic terms.\\
\code{PRAAx} & Defining equations, successor axioms, and quantifier-free
induction.\\
\code{PRA.precZero}, \code{PRA.precSucc} & Rewriting equations for a concrete
primitive-recursive trace runner.\\
\code{PRA.qfInd} & The main invariant rule, with an explicit \code{QF} proof.\\
\code{ConciseIPRA/Native.lean}: \code{NFm}, \code{NPrf} & Native scalar
formulas and proofs.\\
\code{ConciseIPRA/NativeBiInterpretation.lean}: \code{pra_iff_native} &
Faithful export of a completed \code{PRAAx} derivation.\\
\code{ConciseIPRA/NativePublic.lean}: \code{PGPrf}, public gating & Optional
compact, substitution-stable public certificate language.\\
\bottomrule
\end{longtable}

The core arithmetic formalization should target \code{PRAAx}.  This keeps
rational coding and Boolean invariants readable.  Native translation is the
last step.

\subsection{The trace-system data}

After pairing all numerical parameters into one code \code{u}, a Lean
meta-level record can hold the fixed object-language symbols:

\begin{lstlisting}[style=leanstyle]
structure PRTraceSystem where
  good   : PR 1
  init   : PR 1
  bound  : PR 1
  step   : PR 3   -- arguments: i, previousState, u
  inv    : PR 3   -- arguments: i, state, u
  accept : PR 2   -- arguments: state, u

def PRTraceSystem.run (T : PRTraceSystem) : PR 2 :=
  .prec T.init T.step
\end{lstlisting}

The order matches the repository's primitive-recursion convention: the
recursion argument comes first, and the step receives the predecessor, prior
value, and parameters.  Pairing can reduce larger arities to this interface.

Define an object-language predicate \code{IsOne t} as the equation of a
Boolean term with the chosen code for one.  For a concrete \code{T}, construct
terms
\[
 \mathsf{Run}(i,u),\quad \mathsf{Good}(u),\quad
 \mathsf{Inv}(i,u),\quad \mathsf{Accept}(u).
\]
Do not create a second Boolean or rational encoding if the repository already
has one suitable for reuse.

\subsection{Embedded proof of trace reflection}

The dependency-ordered Lean proof is:

\begin{enumerate}
\item Prove simplification lemmas for \code{T.run} by instantiating
      \code{PRA.precZero} and \code{PRA.precSucc}.
\item Define the formula
\[
 \phi(u,i)=\mathsf{Good}(u)\to\mathsf{Inv}(i,u).
\]
      Prove \code{phi.QF} structurally.  Every bounded assertion inside
      \code{Inv} must already be a Boolean-fold equation.
\item Supply the base derivation from the concrete module's base obligation.
\item Under \code{phi :: Gamma.map Fm.shift}, use the induction hypothesis,
      the concrete step obligation, and the successor equation for
      \code{T.run} to derive \code{phi} at successor.
\item Invoke \code{PRA.qfInd phi hQF h0 hS}.
\item Eliminate the universal at \code{T.bound u}, apply final soundness, and
      introduce \code{u} universally.
\end{enumerate}

The generic theorem should take the base, step, and final derivations as fields
of a record.  This makes theorem-specific files provide only local arithmetic
lemmas.

\subsection{Boolean folds and bounded search}

The first supporting module should implement object-language symbols and
proofs for:
\begin{itemize}
\item pair/list projections and length;
\item Boolean \code{all}, \code{any}, and implication;
\item finite sum, product, maximum, and argmin/argmax;
\item first successful bounded search;
\item table lookup with a total default.
\end{itemize}

Each correctness theorem should be a \code{Prf PRAAx} object.  In particular,
\code{argmin} needs the invariant
\[
 \mathsf{bestIndex}<i+1\quad\wedge\quad
 \forall j<i+1\;v(\mathsf{bestIndex})\le v(j),
\]
where the bounded universal part is represented by a Boolean fold.  A
least-index tie rule makes the output deterministic.

\subsection{Finite query tables for names}

For a fixed evaluator symbol \code{a}, define formulas
\code{RegularA}, \code{EqualAB}, and \code{ApartAB}.  A theorem-specific query
planner produces a finite list of precision indices.  The object-language
\code{GoodTable} contains:

\begin{itemize}
\item every instantiated regularity inequality;
\item every domain and magnitude bound;
\item every finite relation between repeated samples;
\item all externally supplied certificate checks.
\end{itemize}

The proof of \code{RegularA -> GoodTable = 1} is separate from trace induction.
Universal elimination is applied at each generated index; a fold combines the
facts.  For adaptive queries, an induction on the trace stage derives the next
instance at the term computed by the current state.

\subsection{Rational and interval kernel}

A canonical object-language rational module should provide:

\begin{enumerate}
\item signed integers and positive-denominator rationals;
\item cross-multiplied equality and order;
\item exact addition, multiplication, absolute value, min/max, powers;
\item dyadics and a function finding a dyadic below a positive rational;
\item interval addition, multiplication with magnitude bounds, and guarded
      reciprocal;
\item finite polynomial evaluation and the Lipschitz majorant of
      \cref{lem:poly-lip}.
\end{enumerate}

The key target is not an abstract ordered-field hierarchy.  It is a small set
of concrete derivations used by the trace checker.  Ambient algebraic
hierarchies may help validate formulas, but the endpoint must use coded
natural-number operations.

\subsection{Robust branch and finite-net modules}

The overlapping branch theorem is a short embedded rational inequality proof.
Its implementation inputs codes \(a,e,q,d\), checks \(4e\le d\), compares
\(a\) with \(q+d/2\), and returns a branch bit plus one of two checked margin
inequalities.  The exact-gap theorem consumes an apartness certificate.

The finite-net module needs:
\begin{itemize}
\item grid generation and a finite coverage theorem for each standard domain;
\item bounded argmin correctness;
\item a Lipschitz transfer lemma;
\item nested-ball regularity and uniqueness.
\end{itemize}
Coverage is domain-specific.  Boxes and simplices should be implemented first;
sphere/projective charts can then reuse robust branch and normalization.

\subsection{Height and complete-list module}

Formalize \cref{thm:height-list} generically over fixed symbols for the point
predicate, height, and bounded enumerator.  The forward direction is a bounded
filter theorem.  The reverse direction is maximum-on-list correctness plus
list completeness.  This module provides a precise reusable endpoint for
Faltings, Siegel, integral-point, and finite-generation packets: they need only
supply the theorem-specific height or covering certificate.

\subsection{Certificate mode and public formulas}

A finite certificate language can be checked in ordinary \code{PRAAx} and
then translated.  If a smaller external proof format is desired, the guarded
public calculus in \code{NativePublic.lean} is relevant: equations are wrapped
by a gating, transport sources are protected, and the public translation is
interderivable with the native kernel on the intended language.  A future
proof compiler can emit public certificates while the trusted Lean theorem
proves their equivalence to \code{NPrf}.

This is an optimization, not the first milestone.  The initial implementation
should favor readable ordinary arithmetic derivations.

\subsection{Suggested module graph}

\begin{center}
\begin{tabular}{@{}l@{}}
\code{CFP/Codes.lean}\\
\quad$\downarrow$\\
\code{CFP/BooleanFold.lean}, \code{CFP/Rational.lean}\\
\quad$\downarrow$\\
\code{CFP/TraceReflection.lean}\\
\quad$\downarrow$\\
\code{CFP/RegularNames.lean}, \code{CFP/RobustBranch.lean}\\
\quad$\downarrow$\\
\code{CFP/FiniteNet.lean}, \code{CFP/PolynomialCertificate.lean},
\code{CFP/HeightSearch.lean}\\
\quad$\downarrow$\\
\text{domain packages and theorem-specific compilers}\\
\quad$\downarrow$\\
\code{pra_iff_native} and optional public-certificate export.
\end{tabular}
\end{center}

\subsection{Milestones and audits}

\begin{enumerate}
\item Compile the finite-trace semantic prototype and prove it without
      \code{sorry}.  This checks the abstract recurrence only.
\item Implement and compile the object-language Boolean-fold library.
\item Construct the first generic \code{PRA.traceReflection} proof.
\item Use it for the small recurrence and robust-branch examples in
      \cref{sec:examples}; inspect \code{\#print axioms}.
\item Formalize regular-name addition, multiplication, and reciprocal under
      apartness.  These are the first all-precision tests of the query-table
      discipline.
\item Formalize polynomial certificate reflection and finite-grid argmin.
\item Port one existing packet's rational core, preferably Euler--Rodrigues or
      finite probability, and replace theorem-specific induction plumbing by
      the generic modules.
\item Add height/list equivalence and port one explicit Diophantine family.
\item Only then optimize certificate size or target the public guarded
      calculus.
\end{enumerate}

Every milestone should run:
\begin{itemize}
\item a source scan for \code{sorry}, \code{admit}, and new line-leading
      \code{axiom};
\item \code{\#print axioms} on exported endpoints;
\item a clean build against the pinned Lean/Mathlib versions;
\item exact regression tests for independent semantic implementations;
\item a status ledger distinguishing semantic, architectural, and embedded
      completion.
\end{itemize}

\subsection{Artifacts supplied in this archive}

\code{lean/FiniteTraceSemantic.lean} gives an ordinary Lean prototype of the
trace invariant and complete-list-to-height argument.  It is not an embedded
proof.  \code{lean/IPRAObjectLanguageBlueprint.lean} imports the actual
\code{ConciseIPRA} interfaces, defines the proposed fixed-symbol record, and
states the completion path.  The nontrivial embedded theorem is deliberately
not claimed complete.  \code{lean/LEAN_PLAN.md} gives the dependency order and
target theorem names.  Lean was unavailable in the artifact-building
environment, so these files were not compiled.
